\chapter{Radial projections and the slope dichotomy}
\label{chap:radial}

We give the numbered ingredients of WZ1 Section 8. Its principal
external input is the thin-tube improvement of Orponen, Shmerkin,
and Wang~\cite{OSW}. WZ1 calls it Lemma 2.9 of that paper; in the
available September 2022 arXiv manuscript the corresponding bootstrap
is Lemma 2.8. WZ1 Proposition 8.6 records the precise version used
here.

\section{Uniform graphs and Frostman sets}

\begin{definition}[Frostman sets; WZ1, Definition 8.1]
\label{def:wz1-frostman}
\uses{def:kt-frostman}
A union $E$ of $\delta$-cubes in $\R^n$ is a
$(\delta,\alpha,C)_n$-Frostman set if
\[
                 |E\cap B(x,r)|\leq Cr^\alpha|E|
                 \quad(r\geq\delta).
\]
For subsets of $S^1$, use unions of $\delta$-arcs. The finite-point
model replaces these volumes by cardinalities of separated
representatives and keeps the normalization by total cardinality.
\end{definition}

\begin{definition}[Fibres and uniform density; WZ1, Definitions 8.2--8.4]
\label{def:uniform-hypergraph}
\lean{Kakeya.Assouad.wz1HypergraphFiber,Kakeya.Assouad.WZ1UniformHypergraphDensity}
\leanok
Let $H\subset A_1\times\cdots\times A_k$. For a coordinate set
$I\subset\{1,\ldots,k\}$ and a projected vertex $a_I$, let
$n_H[a_I]$ be the edges agreeing with $a_I$ on $I$, and let
$\pi_I$ be the coordinate projection. The graph is uniformly
$c$-dense if, for every $h\in H$ and every $I$,
\[
                    \#n_H[\pi_Ih]\geq c\prod_{i\notin I}\#A_i.
\]
The empty and full coordinate sets are included.
\end{definition}

\begin{lemma}[Hypergraph refinement; WZ1, Lemma 8.1]
\label{lem:hypergraph-refinement}
\lean{Kakeya.Assouad.WZ1HypergraphRefinementStatement,Kakeya.Assouad.wz1_hypergraph_refinement}
\leanok
\uses{def:uniform-hypergraph}
If $H$ has density $d=\#H/\prod_i\#A_i$ and $0<\eps<1$,
then it has a uniformly $2^{-k}\eps d$-dense subgraph $H'$ with
$\#H'\geq(1-\eps)\#H$.
\end{lemma}
\begin{proof}
\leanok
For each coordinate projection delete fibres smaller than the
prescribed threshold. Charge the removed edges to their projected
vertices. There are $2^k$ coordinate sets, so the sum of all the
charges is at most $\eps\#H$.
\end{proof}

\begin{theorem}[Robust Kaufman projection; WZ1, Theorem 8.2]
\label{thm:robust-kaufman}
\uses{def:wz1-frostman}
Let $0<\alpha\leq\beta\leq1$ and $\eps>0$.
For sufficiently small $\delta$, suppose $F\subset\R^2$ and
$\Lambda\subset S^1$ are $(\delta,\alpha,C)$- and
$(\delta,\beta,C)$-Frostman sets. Let $H\subset\Lambda\times F$
have density $d$. There is $\Lambda'\subset\Lambda$ with
$|\Lambda'|\gtrsim d|\Lambda|$ such that
\[
                 |\theta\cdot F_\theta|
                       \gtrsim C^{-2}d^3\delta^{1-\alpha+\eps}
                 \quad(\theta\in\Lambda'),
\]
where $F_\theta=\{x:(\theta,x)\in H\}$.
\end{theorem}
\begin{proof}
\uses{thm:kaufman-energy,lem:frostman-energy,lem:energy-frostman}
Discard directions with small fibres or excessive projected energy.
The averaged energy estimate bounds the second exceptional set,
and the density hypothesis controls the first. Cauchy--Schwarz on
each surviving fibre yields the projection lower bound. Choosing
the energy exponent just below $\alpha$ accounts for $\delta^\eps$.
\end{proof}

\section{Thin tubes and improvement of the exponent}

\begin{definition}[Thin tubes; OSW, Definition 2.1; WZ1, Definition 8.5]
\label{def:thin-tubes}
\lean{HasMeasureThinTubes}
\leanok
The ordered pair of probability measures $(\nu_1,\nu_2)$ has
$(\beta,K,1-c)$-thin tubes if there is a Borel relation
$E\subset\supp\nu_1\times\supp\nu_2$ of product measure at
least $1-c$ such that, for every $b_1\in\supp\nu_1$, every line
$\ell$ through $b_1$, and every $r>0$,
\[
          \nu_2\{b_2\in N_r(\ell):(b_1,b_2)\in E\}\leq Kr^\beta.
\]
The reverse ordered pair is a separate condition.
\end{definition}

\begin{definition}[Discrete thin tubes; WZ1, Definition $8.5'$]
\label{def:discrete-thin-tubes}
\lean{Kakeya.Assouad.HasDiscreteThinTubes}
\leanok
For nonempty finite sets $G_1,G_2$, use a relation
$E\subset G_1\times G_2$ with
$\#E\geq(1-c)\#G_1\#G_2$, and require
\[
 \#\{b_2\in G_2\cap N_r(\ell):(b_1,b_2)\in E\}
                                  \leq Kr^\beta\#G_2
                                  \quad(r\geq\delta).
\]
For unions of cells the paper writes the same condition using
Lebesgue measure and a relation that is a union of product cells.
\end{definition}

\begin{lemma}[Returning from smoothed measures; WZ1, Lemma 8.3, comparison]
\label{lem:unsmooth-thin-tubes}
\lean{Kakeya.Assouad.WZ1SmoothedToDiscreteThinTubesStatement}
\uses{def:thin-tubes,def:discrete-thin-tubes}
In the finite-point implementation, smooth each probability measure
at radius $\delta/10$. If these smoothed measures have
$(\beta,K,1-c)$-thin tubes, $0\leq\beta\leq1$ and $2c<1$,
then the original finite sets have discrete
$(\beta,20K,1-2c)$-thin tubes at scale $\delta$.
\end{lemma}
\begin{proof}
Retain a pair of centres when at least half the mass of its product
smoothing cell lies in the good relation. Markov's inequality loses
at most $2c$ of the pairs. Average the measure-level tube estimate
over the source smoothing cell. Moving the centres increases the
tube radius by at most $2\delta/10$, and the half-mass threshold
costs another factor of two; the constant $20$ covers both effects.
\end{proof}

\begin{remark}[The exceptional fraction in Lemma 8.3]
\label{rem:thin-tube-loss}
The paper states a cellwise conversion with constant $4^\beta K/c$
and the same retained fraction $1-c$. Its displayed thresholding
argument, starting with a relation of mass $1-c$ and discarding cells
whose relative good mass is below $c$, bounds the discarded fraction
by $c/(1-c)$, rather than by $c$. The formalized smoothed conversion
uses the explicit loss $2c$ above. Its constants are propagated through
the bootstrap. For uniform measures on an already finite relation,
the direct measure-to-count conversion is a different operation and
does not require this smoothing loss.
\end{remark}

\begin{proposition}[The radial bootstrap; WZ1, Proposition 8.6]
\label{prop:radial-bootstrap}
\lean{radial_bootstrapping_measure_thin_tubes}
\leanok
\uses{def:thin-tubes}
Given $\beta,\eps>0$, there are $\tau=\tau(\beta,\eps)>0$
and $M=M(\beta,\eps)\geq1$ with the following property.
Let $\sigma\in[\beta,1-\eps]$, $0<c<1/10$, and $K,C\geq1$.
Let $\nu_1,\nu_2$ be probability measures supported in the unit
ball of $\R^2$, with support separation at least $1/2$, and
$\nu_i(B(x,r))\leq Cr$ for every positive $r$.
If both ordered pairs have $(\sigma,K,1-c)$-thin tubes, then both
have $(\sigma+\tau,K',1-3c)$-thin tubes, where
\[
                     K'=\max\{K,C^2M/c\}^{M}.
\]
\end{proposition}
\begin{proof}
\uses{thm:os-discretized}
\leanok
If no exponent improvement were possible, choose a scale at which
too many pairs concentrate near lines. The initial thin-tube bounds
make the resulting line families nonconcentrated. A two-ends
alternative separates pairs that spread along their tubes from those
that cluster near one endpoint. In the spread case the discretised
Furstenberg estimate gives too many lines; in the concentrated case
double counting gives a ball with too much measure. The gains are
uniform for $\sigma$ in the prescribed compact interval.
\end{proof}

\begin{remark}[Support normalization in the bootstrap]
\label{rem:radial-support}
The public Lean theorem includes support in the unit ball, whereas
WZ1 Proposition 8.6 writes only the separation and ball-growth
conditions. The original OSW bootstrap is stated in a bounded
normalized region. The application normalizes the two supports
together and tracks the change in their Frostman and thin-tube
constants before invoking it.
\end{remark}

\begin{corollary}[Iterating from exponent $1/4$; WZ1, Corollary 8.7]
\label{cor:thin-tube-iteration}
\uses{def:discrete-thin-tubes}
For every $\eps>0$, a bounded number $N=N(\eps)$ of bootstrap
steps improves discrete $1/4$-thin tubes to $(1-\eps)$-thin tubes
for separated one-dimensional Frostman sets. In the paper's
normalization, initial constants $K,C$ and exceptional fraction $c$
give final exceptional fraction $3^Nc$ and
\[
       K'=\max\{4K/(3c),\,4C^2N/(3c^2)\}^{N}.
\]
At radii $r\geq\delta$, this implies $1$-thin tubes with
constant $\delta^{-\eps}K'$.
\end{corollary}
\begin{proof}
\uses{prop:radial-bootstrap,lem:unsmooth-thin-tubes}
The gain is bounded away from zero while the exponent remains
between $1/4$ and $1-\eps$. Choose the initial exceptional fraction
small enough for every step, and propagate both that fraction and
the growing constant. Finally use
$r^{1-\eps}\leq\delta^{-\eps}r$ for $r\geq\delta$.
In the smoothed formalization the recurrence includes the explicit
conversion losses of Lemma~\ref{lem:unsmooth-thin-tubes}.
\end{proof}

\begin{lemma}[Starting at exponent $1/4$; WZ1, Lemma 8.8]
\label{lem:quarter-thin-tubes}
\lean{Kakeya.Assouad.WZ1QuarterThinTubesStatement,Kakeya.Assouad.wz1_quarter_thin_tubes}
\leanok
\uses{def:wz1-frostman,def:discrete-thin-tubes}
Fix $\lambda,\zeta,\alpha>0$. For sufficiently small $\delta$,
let $G_1,G_2$ be separated $(\delta,1,\delta^{-\lambda})$-Frostman
sets in a bounded region, with mutual distance at least $1/2$.
Assume, for every line $\ell$ and $r\geq\delta$,
\[
            |G_2\cap N_r(\ell)|
                       \leq(\delta^{-\lambda}r)^\zeta|G_2|.
\]
Then $(G_1,G_2)$ has discrete
$(1/4,\delta^{-3\alpha/\zeta-\lambda},1-\delta^\alpha)$-thin
tubes.
\end{lemma}
\begin{proof}
\leanok
Remove pairs lying in excessively populated tubes at a dyadic scale.
The line nonconcentration of $G_2$ and the ball Frostman bounds
control the removed pairs. Sum over the scales and absorb the
logarithm in the chosen power of $\delta$. If the displayed
constant is already large enough for the trivial bound, no removal
is needed.
\end{proof}

\section{Dot differences and separated endpoints}

\begin{definition}[Standard separation; WZ1, Definition 8.6]
\label{def:standard-separation}
\lean{Kakeya.Assouad.WZ1StandardSeparation}
\leanok
The sets $F,G_1,G_2$ have standard separation if each has diameter
at most $1/10$, $\dist(G_1,G_2)\geq1/2$, and
$\dist(F,0)\geq1/2$. They are also placed in a fixed bounded chart.
\end{definition}

\begin{lemma}[Thin tubes give many dot differences; WZ1, Lemma 8.4]
\label{lem:thin-tubes-dots}
\lean{Kakeya.Assouad.WZ1ThinTubesLargeDotProductStatement,Kakeya.Assouad.wz1_thin_tubes_large_dot_product}
\leanok
\uses{def:standard-separation,def:uniform-hypergraph,def:discrete-thin-tubes}
Let $F,G_1,G_2$ be $(\delta,1,K)$-Frostman sets with standard
separation, and assume $(G_1,G_2)$ has discrete
$(1,K,1-c)$-thin tubes. If $H\subset F\times G_1\times G_2$
is uniformly $2c$-dense, then, for sufficiently small $\delta$,
\[
 \cE_\delta\{a\cdot(b_1-b_2):(a,b_1,b_2)\in H\}
                       \gtrsim_\eps c^5K^{-2}\delta^{\eps-1}.
\]
\end{lemma}
\begin{proof}
\uses{thm:robust-kaufman}
\leanok
The thin-tube relation ensures that many endpoint pairs give a
nonconcentrated set of directions $b_1-b_2$. Uniform graph density
keeps substantial coefficient fibres over these pairs. Apply the
robust Kaufman estimate, or its discrete energy proof, and use
standard separation to compare directions with actual differences.
\end{proof}

\begin{proposition}[Line nonconcentration gives a large projection; WZ1, Proposition 8.5]
\label{prop:line-nonconcentration}
\lean{Kakeya.Assouad.wz1_line_nonconcentration_projection}
\uses{def:standard-separation,def:uniform-hypergraph}
For $0<\eps,\zeta<1$, there are $\lambda=\lambda(\eps)>0$
and $\alpha=\alpha(\eps,\zeta)>0$ such that the following holds
at sufficiently small $\delta$. Let $F,G_1,G_2$ be
$(\delta,1,\delta^{-\lambda})$-Frostman sets with standard
separation, and suppose both endpoint sets satisfy
\[
              |G_i\cap N_r(\ell)|
                    \leq(\delta^{-\lambda}r)^\zeta|G_i|
                    \quad(r\geq\delta).
\]
For every uniformly $\delta^\alpha$-dense graph
$H\subset F\times G_1\times G_2$,
\[
       \cE_\delta\{a\cdot(b_1-b_2):(a,b_1,b_2)\in H\}
                                             \geq\delta^{\eps-1}.
\]
\end{proposition}
\begin{proof}
\uses{lem:quarter-thin-tubes,cor:thin-tube-iteration,lem:thin-tubes-dots}
Start with $1/4$-thin tubes in both orders, bootstrap close to
exponent one, and convert to a discrete relation with sufficiently
small exceptional fraction. Choose its fraction below the graph
density, then apply the dot-difference estimate. The choices of
$\lambda$ and $\alpha$ absorb all the constant and exponent losses.
\end{proof}

\begin{proposition}[The separated projection alternative; WZ1, Proposition 8.9]
\label{prop:wz1-separated}
\lean{Kakeya.Assouad.WZ1WellSeparatedProjectionUnionConclusion,Kakeya.Assouad.wz1_well_separated_projection_union}
\leanok
\uses{def:standard-separation,def:uniform-hypergraph}
For every $\eps>0$, sufficiently small $\eta$ and $\delta$ have
the following property. Let $F,G_1,G_2$ be
$(\delta,1,\delta^{-\eta})$-Frostman sets with standard
separation, and let $H\subset F\times G_1\times G_2$ be uniformly
$\delta^\eta$-dense. In the formalized form used for the common
endpoint application, either an affine line $\ell$ and its
perpendicular line $\ell^\perp$ through zero satisfy
\[
 |F\cap N_\delta(\ell^\perp)|_\delta\geq\delta^{\eps-1},
 \qquad
 |G_i\cap N_\delta(\ell)|_\delta\geq\delta^{\eps-1}
 \quad\text{for at least one }i\in\{1,2\},
\]
or the dot-difference set has the long-interval covering lower
bound in Theorem~\ref{thm:wz1-projection}.
\end{proposition}
\begin{proof}
\uses{lem:two-ends-strip,lem:anisotropic-frostman,prop:line-nonconcentration,
  lem:strip-localization,lem:hypergraph-refinement}
\leanok
Apply the two-ends reduction to the endpoint sets, restoring uniform
graph density after each restriction. If the selected strips are
wide, rescale and invoke the line-nonconcentration projection result.
If a strip is narrow, use the localization alternative. The resulting
affine-strip concentration is transported to an active endpoint
class, while the other alternative gives the required long projection.
\end{proof}

\begin{remark}[One endpoint versus two]
\label{rem:endpoint-union}
The printed Proposition 8.9 asks for line concentration in both
$G_1$ and $G_2$ and uses the orientation discussed in
Remark~\ref{rem:projection-orientation}. The closed Lean theorem
linked here returns concentration in at least one endpoint class,
together with concentration of $F$ in the orthogonal line through
zero. This is sufficient in Lemma 5.3, where the endpoint classes
come from one ambient graph. It is a weaker conclusion for
arbitrary independent endpoint sets and is recorded as such.
\end{remark}

\begin{lemma}[Two-ends restriction to a strip; WZ1, Lemma 8.10]
\label{lem:two-ends-strip}
\lean{Kakeya.Assouad.two_ends_reduction}
Let $E\subset[0,1]^2$ be a nonempty union of $\delta$-squares,
and $0<\zeta\leq1/4$. There are a width $w>0$ and a line
$\ell$ such that $E'=E\cap N_w(\ell)$ has
$|E'|\geq\tfrac12 w^\zeta|E|$, while
\[
                    |E'\cap N_r(\ell')|
                       \leq(r/w)^\zeta|E'|
\]
for all lines $\ell'$ and $0<r\leq1$, with harmless fixed
comparison constants in the discrete model.
\end{lemma}
\begin{proof}
Choose a strip nearly maximizing $w^{-\zeta}|E\cap N_w(\ell)|$.
If a narrower strip captured too much of the retained set, it
would contradict this near-maximality. The factor $1/2$ allows
selection without assuming that the supremum is attained.
\end{proof}

\begin{lemma}[Anisotropic Frostman rescaling; WZ1, Lemma 8.11]
\label{lem:anisotropic-frostman}
\lean{Kakeya.Assouad.wz1_anisotropic_frostman_rescaling}
\uses{def:wz1-frostman}
For sufficiently small $\delta$, let
$w\in[\delta^{1-\eps},1]$, and let $E$ be a
$(\delta,1,C)$-Frostman set inside a $1\times w$ rectangle $R$.
If $\phi$ maps $R$ affinely to the unit square, there is
$E'\subset E$ retaining at least $(\delta/w)^\eps$ of its
measure such that the occupied $\delta/w$-squares of $\phi(E')$
form a $(\delta/w,1,(w/\delta)^\eps C)$-Frostman set.
Their intersections with $\phi(E')$ have comparable measure.
\end{lemma}
\begin{proof}
Pigeonhole the masses of the transformed fine cells. Bound the
preimage of a ball by a controlled collection of original-scale
balls, and apply the Frostman condition before normalizing by
the retained mass.
\end{proof}

\begin{lemma}[Frostman coarsening; WZ1, Lemma 8.12]
\label{lem:frostman-coarsening}
\lean{Kakeya.Assouad.wz1_frostman_coarsening}
\uses{def:wz1-frostman}
For $\delta\leq\rho\leq1$, a
$(\delta,\alpha,C)$-Frostman set has a subset retaining a
fraction $\gtrsim(1+\log(1/\delta))^{-1}$ whose occupied
$\rho$-cells form a $(\rho,\alpha,C')$-Frostman set, with
$C'\lesssim C(1+\log(1/\delta))$.
\end{lemma}
\begin{proof}
Group occupied $\rho$-cells by their fine mass and retain a
populous dyadic class. Comparable cell masses let the fine
Frostman estimate be converted into a coarse cell count.
\end{proof}

\begin{lemma}[Localization inside a strip; WZ1, Lemma 8.13]
\label{lem:strip-localization}
\lean{Kakeya.Assouad.WZ1StripLocalizationDichotomyStatement}
\uses{def:standard-separation,def:uniform-hypergraph}
For every $\eps>0$, choose $\eta$ and $\delta$ sufficiently
small. Suppose $F,G_1,G_2$ are
$(\delta,1,\delta^{-\eta})$-Frostman sets with standard separation,
$H$ is uniformly $\delta^\eta$-dense, and
$G_1\subset N_w(\ell)$, $w\geq\delta$. Then either
\[
 \pi_{G_2}H\subset N_{\delta^{-\eps}w}(\ell),\qquad
 \pi_FH\subset N_{\delta^{-\eps}w}(\ell^\perp),
\]
where $\ell^\perp$ passes through zero, or the dot-difference set
has the long-interval lower bound of Theorem~\ref{thm:wz1-projection}.
\end{lemma}
\begin{proof}
\uses{thm:robust-kaufman,lem:anisotropic-frostman,lem:frostman-coarsening}
If an active endpoint or coefficient lies outside the proposed
strips, its separation yields a nondegenerate projection scale.
Rescale the common active strip, coarsen with controlled Frostman
loss, and apply Kaufman's projection estimate to the resulting
directions. Transport that lower bound back to the original
dot-difference graph.
\end{proof}

\begin{remark}[Active vertices]
\label{rem:active-vertices}
The first alternative concerns vertices occurring in $H$. Unused
members of the ambient $F,G_1,G_2$ need not lie in either strip.
The code measures the common strip width on the active coordinate
projections and truncates it below by $\delta$. This matches the
surjective-projection reduction in the paper and prevents unrelated
ambient points from changing the rescaling factor.
\end{remark}
